\section{System Model}
\label{sec:system-model}

\subsection{Architecture Context}

The Basis Network zkEVM L2 employs a per-enterprise chain model
where each enterprise operates its own sequencer, prover, and state
database. The state database serves as the authoritative record of
the EVM world state, mapping account addresses to their
current state (nonce, balance, code hash, storage root). The
sequencer updates the state database after executing each
transaction, and the prover reads Merkle witnesses from it to
construct validity proofs that are verified on the Basis Network L1.

The state database occupies a critical position in the block
production pipeline:

\begin{enumerate}
  \item Sequencer receives transactions and orders them into a block.
  \item EVM executor processes each transaction, producing state
    changes (account updates, storage modifications).
  \item \textbf{State database applies updates and computes the new
    state root.}
  \item The state root is included in the block header.
  \item The prover reads Merkle witnesses from the state database
    to generate a ZK proof of the state transition.
\end{enumerate}

Step 3 is the subject of this paper. Its latency directly bounds
block production throughput: if state root computation exceeds the
block time budget, the sequencer cannot produce blocks at the target
rate.

\subsection{State Tree Requirements}

The state database must satisfy the following requirements:

\begin{table}[h]
\centering
\caption{State database requirements for Basis Network zkEVM L2}
\label{tab:requirements}
\begin{tabular}{lll}
\toprule
\textbf{Requirement} & \textbf{Target} & \textbf{Rationale} \\
\midrule
State root latency & $< 50$~ms/block & Sub-second finality \\
Account capacity & $\geq 10{,}000$ & Enterprise scale \\
Memory per account & $< 100$~KB & Commodity hardware \\
Proof generation & $< 1$~ms/proof & Witness extraction \\
Proof verification & $< 1$~ms/proof & Local consistency \\
ZK compatibility & BN254 field & Circuit integration \\
\bottomrule
\end{tabular}
\end{table}

The 50~ms budget derives from the target block time of 500~ms,
allocating 10\% to state root computation. The remaining budget
covers transaction execution ($\sim$200~ms), consensus
($\sim$100~ms), and network propagation ($\sim$150~ms).

\subsection{Data Model}

Each EVM account maps to an SMT leaf:

\begin{equation}
\text{leaf}(a) = H_{\text{Poseidon2}}(\text{key}(a), \text{value}(a))
\label{eq:leaf}
\end{equation}

where $\text{key}(a)$ is derived from the account address and
$\text{value}(a)$ encodes the account state. The state root
$R$ is the root hash of the SMT containing all account leaves:

\begin{equation}
R = \text{SMT}_{\text{root}}(\{\text{leaf}(a) : a \in \text{Accounts}\})
\label{eq:root}
\end{equation}

For a block containing $n$ transactions that modify accounts
$\{a_1, \ldots, a_k\}$ (where $k \leq n$), the state root update
requires:

\begin{equation}
\text{Cost}_{\text{update}} = k \cdot d \cdot C_{\text{hash}}
\label{eq:cost}
\end{equation}

where $d$ is the tree depth and $C_{\text{hash}}$ is the cost of
one Poseidon2 2-to-1 hash compression. This linear relationship
between batch size, depth, and hash cost defines the optimization
space.

\subsection{Threat Model}

The state database operates within an enterprise-controlled
environment where the sequencer is trusted but the state root
must be independently verifiable via ZK proofs on L1. The primary
threats are:

\begin{itemize}
  \item \textbf{State divergence:} Inconsistency between the
    committed state root and the actual account state, detectable
    by the ZK prover.
  \item \textbf{Proof forgery:} Fabricated Merkle proofs for
    non-existent state, prevented by the collision resistance
    of Poseidon2.
  \item \textbf{Performance degradation:} Pathological key
    distributions causing worst-case tree behavior, mitigated
    by key hashing.
\end{itemize}
